% HOUSE RULE: a node gets \lean{} and \leanok in the same commit as the declaration it
% names, never earlier.  The dependency graph is this project's only record of progress,
% so a tag that runs ahead of the code makes it lie.
%
% STATUS (2026-09-19, after Q20).  Everything in RBM3D/ compiles (`./check.sh`), with no
% `sorry` and -- since Q19 -- no project axiom at all.  The nodes below therefore carry
% \lean{} tags, and \leanok on every node that is *proved*.
%
% A node that names an unproved premise gets \lean{} but NOT \leanok: the statement
% exists in Lean, the proof does not.  These are hypotheses carried by the theorems that
% use them -- not axioms -- so the audit reports 0 axioms while still reporting how many
% theorems rest on each of them.  The label prefix says WHY the proof is missing, and it
% matches the two ledgers of RBM3D/Test/Axioms.lean exactly:
%
%   bor:  BORROWED -- the paper cites it rather than proving it.  Discharging one means
%         proving something the paper itself does not.
%         bor:ThfadC, bor:ThfadCshort, bor:BD1, bor:BD2, bor:ThfadC0  (the five
%         components of RBM.PropTH) and bor:KTreeRep ([YY_25] Lemma 3.4).
%   hyp:  OWED -- the paper does prove it; it is assumed here only to get on with the
%         next layer.  A debt of this formalization, not of the paper.
%         hyp:Kbound (proved in Appendix A.5), hyp:twoloop.
%
% Anything else with an unproved statement is a definition of the objects under study
% (def:kloop's IsKLoop), which the audit keeps in a third, structural list.

\chapter{The lattice, the model and the propagator}

% RBM.Zd, RBM.zdistD
\begin{definition}[The torus and its periodic distance]\label{def:lattice}
  \lean{RBM.Zd, RBM.zdistD}\leanok
  Fix $d\ge3$ and $L\ge3$ and work on $\ZL$.  For $x\in\ZL$ let $|x|$ be the periodic
  $\ell^1$ distance to the origin: coordinatewise graph distance on the cycle
  $\ZLone$, summed over the $d$ coordinates.  The paper does not fix a norm and notes
  in \S2.1 that the choice is immaterial; see \texttt{docs/paper-deltas.md}, D2.
\end{definition}

% RBM.SB, (eq:variancematrix)
\begin{definition}[Block variance matrix $\SB(\gcoup)$]\label{def:SB}
  \lean{RBM.SB}\leanok
  \uses{def:lattice}
  For a coupling parameter $\gcoup>0$,
  \[ \SB_{ab}(\gcoup) \;=\; \frac{1}{1+2d\gcoup^2}\,\mathbf 1_{a=b}
     \;+\; \frac{\gcoup^2}{1+2d\gcoup^2}\,\mathbf 1_{a\sim b},\qquad a,b\in\ZL, \]
  where $a\sim b$ means $|a-b|=1$.  Since the right-hand side depends only on $a-b$,
  $\SB(\gcoup)$ is the circulant matrix on the group $\ZL$ generated by that kernel.
\end{definition}

\begin{remark}[Why $L\ge3$]\label{rem:L3}
  The paper leaves $L\ge3$ implicit.  It is needed: for $L\le2$ the $2d$ neighbours of a
  point are not distinct --- for $L=2$ one has $1=-1$ in $\ZLone$ --- and $\SB(\gcoup)$
  fails to be doubly stochastic.  See \texttt{docs/paper-deltas.md}, D1.
\end{remark}

% RBM.card_nbhd, RBM.card_adj  (RBM3D/Defs/Neighbours.lean)
\begin{lemma}[Neighbour count]\label{lem:card-nbhd}
  \lean{RBM.card_nbhd, RBM.card_adj}\leanok
  \uses{def:lattice}
  For $L\ge3$, $\#\{x\in\ZL : |x|=1\}=2d$, and every $a\in\ZL$ has exactly $2d$
  neighbours.  The unit sphere is the image of the injection $(i,\pm)\mapsto\pm e_i$ of
  $\{1,\dots,d\}\times\{\pm\}$; injectivity is $1\ne-1$ in $\ZLone$, which needs $L\ge3$.
\end{lemma}

% RBM.SB_isSymm, RBM.SB_apply_add_right, RBM.sum_SB_row, RBM.SB_mulVec_one, RBM.norm_SB
\begin{lemma}[Elementary properties of $\SB$]\label{lem:SB-basic}
  \lean{RBM.SB_isSymm, RBM.SB_apply_add_right, RBM.sum_SB_row, RBM.SB_mulVec_one, RBM.norm_SB}\leanok
  \uses{def:SB, lem:card-nbhd}
  $\SB(\gcoup)$ is symmetric and translation invariant, each row sums to $1$, and
  $\SB(\gcoup)\mathbf 1=\mathbf 1$.  The row sum is where $L\ge3$ enters: it amounts to
  the count $\#\{x\in\ZL : |x|=1\}=2d$.
\end{lemma}

% RBM.ellT (eq:ellt), RBM.Bparam (eq_B_param)
\begin{definition}[Control parameters $\ell_t$ and $B_{t,K}$]\label{def:params}
  \lean{RBM.ellT, RBM.Bparam}\leanok
  \uses{def:lattice}
  For $t\in[0,1)$,
  \[ \ell_t \;=\; \min\bigl(\max(\gcoup|1-t|^{-1/2},1),\,L\bigr),\qquad
     B_{t,K} \;=\; \frac{(\gcoup^2+|1-t|)^{-1}}{(K+1)^{d-2}} + \frac{1}{L^d|1-t|}. \]
  $\ell_t$ is the effective range of the propagator and $B_{t,K}$ its decay profile.
\end{definition}

% RBM.Theta, RBM.ThetaRBM (def_Thxi), RBM.ThetaRBM0 (def_Thxi0)
\begin{definition}[The $\Th$-propagator]\label{def:Theta}
  \lean{RBM.Theta, RBM.Theta0}\leanok
  \uses{def:SB}
  For $t\in[0,1]$ and $\sigma_1,\sigma_2\in\{+,-\}$,
  \[ \Th_t^{(\sigma_1,\sigma_2)} \;=\; \bigl(1 - t\,M^{(\sigma_1,\sigma_2)}\SB\bigr)^{-1},
     \qquad
     \mathring\Th_t(a,b) \;=\; \Th_t(a,b) - L^{-2d}\!\!\sum_{a',b'}\!\Th_t(a',b'). \]
  For the random band matrix model $M^{(\sigma_1,\sigma_2)}=m(\sigma_1)m(\sigma_2)I$, so
  the propagator depends on the two signs only through the product
  $m=m(\sigma_1)m(\sigma_2)$; that is the case formalized.  See
  \texttt{docs/paper-deltas.md}, D6.
\end{definition}

% RBM.UnifDetDom, RBM.DetDom
\begin{definition}[Deterministic stochastic domination]\label{def:prec}
  \lean{RBM.UnifDetDom, RBM.DetDom}\leanok
  $\xi\prec\zeta$ if for every $\tau>0$ one has $\xi\le N^\tau\zeta$ for all large $N$,
  uniformly in any parameters present.
\end{definition}

\chapter{Properties of the propagator: what the paper proves}

% RBM.Theta_isSymm  (T2)
\begin{lemma}[$\mathtt{lem\_propTH}$, properties 1--2]\label{lem:propTH-12}
  \lean{RBM.Theta_isSymm}\leanok
  \uses{def:Theta, lem:SB-basic}
  $\Th_t^{(\sigma_2,\sigma_1)}=(\Th_t^{(\sigma_1,\sigma_2)})^\top=\Th_t^{(\sigma_1,\sigma_2)}$,
  and $\Th_t(a+r,b+r)=\Th_t(a,b)$ for all $a,b,r\in\ZL$.
  Appendix A.1 proves both from the corresponding properties of $M^{(\sigma_1,\sigma_2)}\SB$
  through the Taylor expansion $\Th_t=\sum_k t^k(M\SB)^k$.
\end{lemma}

% RBM.Theta_commute_SB_of_three_le, RBM.Theta_commute_of_three_le, RBM.norm_Theta_apply_le, RBM.sum_Theta_real_row, RBM.norm_Theta_le  (Propagator/Props4.lean)
\begin{lemma}[$\mathtt{lem\_propTH}$, properties 3--4]\label{lem:propTH-34}
  \lean{RBM.Theta_commute_SB_of_three_le, RBM.Theta_commute_of_three_le, RBM.norm_Theta_apply_le, RBM.sum_Theta_real_row, RBM.norm_Theta_le}\leanok
  \uses{def:Theta, lem:SB-basic}
  $[\SB,\Th_t]=[\Th_t,\Th_{t'}]=0$; and $|\Th_{t,ab}^{(\sigma_1,\sigma_2)}|\le\Th_{t,ab}^{(+,-)}$
  with
  \[ \|\Th_t^{(\sigma_1,\sigma_2)}\|_{\infty\to\infty}
     \;=\;\max_a\sum_b|\Th_{t,ab}^{(\sigma_1,\sigma_2)}|\;\le\;\frac{1}{1-t}. \]
  Appendix A.1 proves this in full: termwise comparison in the Taylor expansion, then
  $\sum_b\Th_{t,ab}^{(+,-)}=(1-t)^{-1}$ because $M^{(+,-)}\SB$ is doubly stochastic.
\end{lemma}

\chapter{What is assumed, and why: the two ledgers}

Nothing in this development is an axiom (Q19: the audit
\texttt{RBM3D/Test/Axioms.lean} reports 0, and fails the build if an unregistered premise
appears).  What is unproved is unproved \emph{as a hypothesis}: a \texttt{Prop} carried by
the theorems that need it.  Those premises fall into two kinds, and the audit keeps two
separate ledgers for them, because reporting them in one list makes the situation look
better than it is.  The label prefixes below say which ledger a node is in.

\paragraph{Borrowed (\texttt{bor:}) --- the paper cites it rather than proving it.}
Five of them are the components of $\mathtt{lem\_propTH}$ collected in this chapter,
bundled in Lean as the structure \texttt{RBM.PropTH}
(\texttt{RBM3D/Propagator/Interface.lean}); the sixth is the tree representation
\ref{bor:KTreeRep}, \texttt{[YY\_25]} Lemma 3.4.  Discharging one of these means proving
something \emph{this paper does not prove}, so they are long-term: the list is the precise,
machine-checkable record of what the paper borrows from the literature.  Q21 checked the
five statements here against the paper word by word: one was \emph{false} as first written
and had to be repaired (see \texttt{docs/paper-deltas.md}, D11, and the machine-checked
counterexample \texttt{RBM.not\_decayShort\_at\_one}), and one was weaker than the paper's
claim (D12).

\paragraph{Owed (\texttt{hyp:}) --- the paper proves it; this formalization does not, yet.}
Two of them: \ref{hyp:Kbound}, whose proof Appendix A.5 gives in full (``for the reader's
convenience, we provide the proof below'') but which needs the molecule layer, and
\ref{hyp:twoloop}, the a priori bound on $2$-loops.  These are debts of the Lean
development, and discharging one changes nothing about the paper.  They are the ones to
look at first: a shrinking \texttt{hyp:} list is progress, a shrinking \texttt{bor:} list
is a new theorem.

\begin{theorem}[$\mathtt{(prop{:}ThfadC)}$, property 5 --- \emph{interface hypothesis}]\label{bor:ThfadC}
  \lean{RBM.ThetaDecay}
  \uses{def:Theta, def:params}
  There are $C_d,c_d>0$ with
  $|\Th_t(0,a)|\le C_d\,B_{t,|a|}\,e^{-c_d|a|/\ell_t}$ for all $a\in\ZL$.
  \emph{Attributed to} \texttt{[DYYY25]} Lemma 2.14; Appendix A.1 sketches the
  random-walk argument and then says ``we omit the details''.
\end{theorem}

\begin{theorem}[$\mathtt{(prop{:}ThfadC\_short)}$ --- \emph{interface hypothesis}]\label{bor:ThfadCshort}
  \lean{RBM.ThetaDecayShort}
  \uses{def:Theta}
  For $\sigma_1=\sigma_2$ there are $C_\kappa,c_\kappa>0$ with
  $|\Th_t(0,a)|\le C_\kappa(\mathbf 1_{a=0}+\gcoup^2e^{-c_\kappa|a|})$.
  \emph{Proved} in Appendix A.1, but through \texttt{[bourgade2019random]} Lemma 4.2,
  which is outside this paper.
\end{theorem}

\begin{theorem}[$\mathtt{(prop{:}BD1)}$, property 6 --- \emph{interface hypothesis}]\label{bor:BD1}
  \lean{RBM.ThetaDiffOne}
  \uses{def:Theta, def:prec}
  For $|r|\lesssim|a|$,
  $|\Th_t(0,a+r)-\Th_t(0,a)|\prec(\gcoup^2+|1-t|)^{-1}|r|(|a|+1)^{-(d-1)}$.
  \emph{Not stated explicitly anywhere the paper points to}: ``Although the bound is not
  stated explicitly in \texttt{[yang2024Del]}, its proof proceeds analogously \dots we
  therefore omit the details.''  This is the one estimate with no proof in the cited
  literature at all.
\end{theorem}

\begin{theorem}[$\mathtt{(prop{:}BD2)}$, property 7 --- \emph{interface hypothesis}]\label{bor:BD2}
  \lean{RBM.ThetaDiffTwo}
  \uses{def:Theta, def:prec}
  For $|r|\lesssim|a|$,
  $|\Th_t(0,a+r)+\Th_t(0,a-r)-2\Th_t(0,a)|\prec(\gcoup^2+|1-t|)^{-1}|r|^2(|a|+1)^{-d}$.
  \emph{Attributed to} \texttt{[yang2024Del]} (E.19).
\end{theorem}

\begin{theorem}[$\mathtt{(prop{:}ThfadC0)}$, property 8 --- \emph{interface hypothesis}]\label{bor:ThfadC0}
  \lean{RBM.ThetaZeroMode}
  \uses{def:Theta, def:prec}
  $|\mathring\Th_t(0,a)|\prec(\gcoup^2+|1-t|)^{-1}(|a|+1)^{-(d-2)}$.
  \emph{Attributed to} \texttt{[yang2024Del]} Lemma 3.1.
\end{theorem}

\begin{lemma}[Certificates: which of these are known to be satisfiable]\label{lem:certificates}
  \lean{RBM.Test.thetaDecay_fixedL, RBM.Test.exists_norm_Theta_ge, RBM.Test.not_exists_uniform_entry_bound, RBM.Test.sum_Theta_diff_row, RBM.Test.twoLoopBounded_kTwoLoop, RBM.Theta0_apply_eq, RBM.sum_Theta0_row}\leanok
  \uses{bor:ThfadC, bor:ThfadCshort, bor:BD1, bor:BD2, bor:ThfadC0, hyp:twoloop, lem:propTH-34}
  The audit counts what rests on each premise, but a premise that is \emph{false} makes
  every theorem carrying it vacuous, and the count looks the same.  That is not
  hypothetical: it is what happened while $\mathtt{(prop{:}ThfadC\_short)}$ was stated
  without its $\sigma_1=\sigma_2$ restriction (\texttt{docs/paper-deltas.md}, D11).  So
  each premise also wants a \emph{certificate}: a theorem of this development witnessing
  that it is satisfiable.  Two have one.

  \ref{bor:ThfadC} is certified at \textbf{fixed $L$}: replacing
  $\exists C\,\forall L$ by $\forall L\,\exists C$ leaves a statement this development can
  prove, because at fixed $L$ the torus is finite and the whole content of the assumption
  is the uniformity in $L$.  The proof plays $\|\Th_t\|_{\infty\to\infty}\le(1-t)^{-1}$
  against the zero-mode term $(L^d|1-t|)^{-1}$ of $B_{t,K}$, and
  \texttt{exists\_norm\_Theta\_ge} shows that term is not slack: some entry of the row is
  at least $L^{-d}(1-t)^{-1}$.  \ref{hyp:twoloop} is certified outright, by the explicit
  two-loop of $\mathtt{(Kn2sol)}$.

  The other four propagator premises are \textbf{not} certified, and the obstruction is
  machine-checked rather than asserted: their right-hand sides stay bounded as $t\to1$ at
  fixed $L$, while the entries of $\Th_t$ do not
  (\texttt{not\_exists\_uniform\_entry\_bound}), so each of them needs a
  \emph{cancellation} --- a difference of entries, or the removal of the zero mode.  At
  fixed $L$ that cancellation is the statement that $1$ is a \emph{simple} eigenvalue of
  $\SB$, a finite-$L$ spectral gap this development does not have (work order Q51).  What
  is checked is that the refutation which killed the old property 5' cannot be rerun
  against them: $\sum_b\mathring\Th_{ab}=0$ and every first difference of a row sums to
  $0$.
\end{lemma}

\chapter{Scaling order (Appendix B)}

% RBM.Graph.Counters, RBM.Graph.ord  (eq:ordG)
\begin{lemma}[The derivative of the propagator, $\mathtt{(2.51)}$ of the $d=1$ paper]\label{lem:Theta-deriv}
  \lean{RBM.continuousAt_Theta, RBM.Theta_sub_Theta, RBM.continuous_matrix_entry, RBM.hasDerivAt_Theta_apply, RBM.hasDerivAt_Theta_mul_apply}\leanok
  \uses{def:Theta, lem:SB-basic}
  Entrywise, $\partial_\xi(\Th_\xi)_{ab}=(\Th_\xi\SB\Th_\xi)_{ab}$ for $\|\xi\|<1$, and
  along $\xi=t\mu$, $\partial_t(\Th_{t\mu})_{ab}=\mu(\Th_{t\mu}\SB\Th_{t\mu})_{ab}$.
  The proof goes through the resolvent identity
  $\Th_\zeta-\Th_\xi=(\zeta-\xi)\Th_\zeta\SB\Th_\xi$ and continuity of $\xi\mapsto\Th_\xi$,
  not by differentiating the Neumann series.

  The statement is \emph{entrywise} on purpose: that is what the loop layer uses, and it
  avoids the instance diamond between the Pi topology on \texttt{Matrix} and the topology
  of the $\ell^\infty$ operator norm.  This node is the prerequisite for \ref{lem:KTwoFormula} (done: Q27 and Q22a), for
  \ref{lem:treeThree} ($n=3$, Q29) and for \ref{bor:KTreeRep}, which now lacks only
  existence at $n\ge4$ (Q30, Q31).
\end{lemma}

\begin{definition}[Scaling size and scaling order]\label{def:ord}
  \lean{RBM.Graph.Counters, RBM.Graph.ord}\leanok
  For a normal graph $\Gamma$ with $n_S$ solid edges (light-weights included), $n_W$ waved
  edges, $n_V$ internal vertices and $n_M$ internal molecules,
  \[ \ord(\Gamma) \;=\; n_S(\Gamma) + 2\bigl(n_W(\Gamma)-n_V(\Gamma)\bigr). \]
  The Appendix B variant $\mathtt{(eq{:}ordG\_BA)}$ replaces internal vertices by internal
  atoms; the formula is the same.  $\ord$ takes values in $\mathbb Z$, since $n_W-n_V$ may
  be negative.  See \texttt{docs/paper-deltas.md}, D7.
\end{definition}

% RBM.Graph.ord_case_ii .. ord_case_vi
\begin{lemma}[Counter bookkeeping in $\mathtt{lem\_scalingorder}$]\label{lem:ord-cases}
  \lean{RBM.Graph.ord_case_ii, RBM.Graph.ord_case_vi}\leanok
  \uses{def:ord}
  In each configuration of a weight expansion, the counter relations between $\Gamma_0$ and
  $\Gamma_1$ force a lower bound on $\ord(\Gamma_1)-\ord(\Gamma_0)$: $+3$ in cases (ii) and
  (v), $+2$ in cases (iii), (iv) and (vi).
\end{lemma}

\begin{remark}[What the arithmetic does \emph{not} capture]\label{rem:ord-exhaustive}
  \uses{lem:ord-cases}
  \Cref{lem:ord-cases} is pure arithmetic.  The mathematical content of the case analysis
  is that the listed configurations are \emph{exhaustive}, and that is a statement about
  the graph model, not about the counters.

  This is not a hypothetical concern.  The third proofreading round of the paper added a
  case (vi) to this analysis that had been omitted --- and its counter relations are
  \emph{subsumed} by those of case (iv), so \texttt{ord\_case\_vi} is a corollary of
  \texttt{ord\_case\_iv}.  A formalization that checked only the arithmetic would not have
  caught the omission.  Building the graph model so that the case split is exhaustive by
  construction (T5) is therefore where this project earns its keep.
\end{remark}

% RBM.Graph.hasDerivAt_inverse_apply, RBM.Graph.hasDerivAt_inverse_sub_apply  (Graph/Expansions.lean)
\begin{lemma}[Derivative of a resolvent entry]\label{lem:dG}
  \lean{RBM.Graph.hasDerivAt_inverse_apply, RBM.Graph.hasDerivAt_inverse_sub_apply}\leanok
  For invertible $A$ with $G=A^{-1}$, $\partial_{h_{\alpha w}}G_{ij}=-G_{i\alpha}G_{wj}$, and the same for
  $\mathring G=G-M$.  In the third term of $\mathtt{(Owx)}$ this turns a solid edge $G_{\beta_1\beta_2}$
  into $G_{\beta_1\alpha}G_{w\beta_2}$; with the prefactor $G_{\alpha w}$ these are the three new
  edges classified in \cref{lem:ord-exhaustive}.  The expansions $\mathtt{(Owx)}$, $\mathtt{(Oe2x)}$
  themselves (\texttt{[yang2021delocalization]} Lemmas 3.5, 3.14) are identities in expectation
  and are not formalized; see \texttt{docs/QUEUE.md}, Q8.
\end{lemma}

% RBM.Graph.Pattern.classify, RBM.Graph.classify_eq_vi_iff, RBM.Graph.ord_weight_step
\begin{lemma}[The case split of $\mathtt{lem\_scalingorder}$ is exhaustive]\label{lem:ord-exhaustive}
  \lean{RBM.Graph.Pattern.classify, RBM.Graph.classify_eq_vi_iff, RBM.Graph.ord_weight_step}\leanok
  \uses{lem:ord-cases, lem:dG}
  A configuration of the weight-expansion step is the equality pattern of the vertices on
  the cycle $\beta_1 - \alpha - w - \beta_2 - \beta_1$.  Of the $16$ Boolean patterns, exactly
  $12$ are realizable (no $4$-cycle has exactly three coincidences), and each of them falls
  into one of cases (i)--(vi); each case is characterised by the vertex (in)equalities the
  paper uses to describe it.  Consequently $\ord(\mathcal G_1) \ge \ord(\mathcal G_0)+1$ in
  every configuration, given the counter relations stated for each case.  In Lean the split
  is a single \texttt{match} whose exhaustiveness the compiler checks: removing case (vi)
  is rejected with \texttt{missing cases}.
\end{lemma}

\chapter{Deterministic estimates (Appendix A)}

The remaining chapters are stated but not yet formalized; they are work orders T6--T8.

% RBM.UN, RBM.ThetaN (DefTHUST); RBM.uKer_eq_one_add, RBM.norm_Xi_le, RBM.norm_UN_le  (Kernel/Evolution.lean)
\begin{lemma}[$\mathtt{lem{:}sum\_Ndecay}$, Appendix A.2]\label{lem:sum-Ndecay}
  \lean{RBM.UN, RBM.ThetaN, RBM.uKer_eq_one_add, RBM.norm_UN_le}\leanok
  \uses{lem:propTH-34}
  For $0\le s\le t<1$,
  $\|\mathcal U^{(n)}_{s,t,\bm\sigma}\circ\mathcal A\|_\infty\le\bigl(\tfrac{1-s}{1-t}\bigr)^n\|\mathcal A\|_\infty$.
  Each one-index factor is $1+\Xi^{(i)}$ with $\|\Xi^{(i)}\|_{\infty\to\infty}\le(t-s)/(1-t)$ by property 4.
\end{lemma}

\begin{lemma}[Evolution kernel decomposition, Appendix A.2]\label{lem:Ukernel}
  \lean{RBM.tensorKer, RBM.projMat, RBM.zeroModeSet, RBM.norm_zeroModeSet_UN_le}\leanok
  \uses{bor:ThfadC, lem:propTH-34}
  The kernel $\Ucal^{(n)}_{s,t,\bm\sigma}$ acts on each of the $n$ indices separately, and
  $Q^{(A)}$ replaces the factors indexed by $i\in A$ by ones carrying
  $\mathsf{Proj}_{\mathbf e^\perp}$.  For $i\notin A$ the assumption
  $A\supset I_{\mathrm{diff}}(\bm\sigma)$ forces $\sigma_i=\sigma_{i+1}$.
\end{lemma}

% RBM.tailT, RBM.tailW, RBM.tailT_antitone, RBM.tailW_pos, RBM.zeroMode_le_of_ge, RBM.ellT_eq_of_le  (Defs/Tail.lean)
\begin{definition}[Tail functions, $\mathtt{def{:}\ TTfunc}$]\label{def:tail}
  \lean{RBM.tailT, RBM.tailW, RBM.tailT_antitone, RBM.tailW_pos, RBM.zeroMode_le_of_ge, RBM.ellT_eq_of_le}\leanok
  \uses{def:params}
  $\mathcal T_t(r)=B_{t,r}\,e^{-(r/\ell_t)^{1/2}}$ and
  $\widetilde{\mathcal T}^{\ell}_{t,D}(r)=\max(\mathcal T_t(r\wedge\ell),W^{-D})$.  $\mathcal T_t$ is
  non-negative and non-increasing on $[0,\infty)$ with $\mathcal T_t(0)=B_{t,0}$;
  $\widetilde{\mathcal T}\ge W^{-D}>0$.  For $0\le r\le L$: if $1-t\ge g^2/L^2$ then
  $(L^d|1-t|)^{-1}\le 2^{d-1}(g^2+|1-t|)^{-1}(r+1)^{-(d-2)}$; if $1-t\le g^2/L^2$ then
  $\ell_t=L$ and $e^{-(r/\ell_t)^{1/2}}\ge e^{-1}$.
\end{definition}

% RBM.propT, RBM.propT_i, RBM.propT_ii  (Kernel/PropT.lean); RBM.sum_conv_le  (Defs/Convolution.lean)
\begin{lemma}[$\mathtt{lem{:}propT}$, $\mathtt{(TTT2)}$]\label{lem:propT-TTT2}
  \lean{RBM.propT, RBM.propT_i, RBM.propT_ii, RBM.sum_conv_le}\leanok
  \uses{def:tail}
  There is $C_d>0$ such that for $0\le u\le t<1$ with (i) $1-u\ge1-t\ge g^2/L^2$ or
  (ii) $1-t\le1-u\le g^2/L^2$, and all $a,b$,
  $\sum_c\mathcal T_u(|a-c|)\mathcal T_t(|c-b|)\le\frac{C_d}{1-u}\mathcal T_t(|a-b|)$.
  The ``basic calculus fact'' on $\mathbb R^d$ used in Appendix A.3 is replaced by its
  lattice form, proved via $\sqrt{p/\ell_1}+\sqrt{q/\ell_2}-\sqrt{r/\ell_2}\ge(2-\sqrt2)\sqrt{\min(p,q)/\ell_1}$ for $r\le p+q$.
\end{lemma}


% ---------------------------------------------------------------------------
% Appendix A.4: the truncated tail function and claim:TTk
% ---------------------------------------------------------------------------

\begin{definition}[The truncated tail $\mathsf T_t$ and $\Psi_t$, Appendix A.4]\label{def:sfT}
  \lean{RBM.sfT, RBM.PsiT, RBM.sfT_zero_le_PsiT, RBM.sfT_le_PsiT_mul, RBM.sfT_antitone, RBM.wfac}\leanok
  \uses{def:tail, def:params}
  $\mathsf T_t(r)=(\gcoup^2+|1-t|)^{-1/2}W^{-d/2}(r+1)^{-(d-2)/2}e^{-\frac12\sqrt{r/\ell_t}}$ and
  $\Psi_t=(W^{-d}B_{t,0})^{1/2}$.  $\mathsf T_t$ is non-negative and non-increasing, and the
  two elementary facts of Appendix A.4 hold: $\mathsf T_t(0)\le\Psi_t$ and
  $\mathsf T_t(r)\le\mathsf T_t(0)(r+1)^{-(d-2)/2}$.  We write $w_d(r)=(r+1)^{-(d-2)/2}$ for the
  polynomial factor that $\mathtt{(eq{:}TtTt)}$ and $\mathtt{(eq{:}KtKt)}$ carry.
\end{definition}

\begin{lemma}[$\mathtt{claim{:}TTk}$ pointwise: $\mathtt{(eq{:}TtTt)}$, $\mathtt{(eq{:}KtKt)}$ and their coverage]\label{lem:TTk-pointwise}
  \lean{RBM.sfT_mul_le_TtTt, RBM.sfT_mul_le_KtKt, RBM.sfT_TtTt, RBM.sfT_KtKt, RBM.sfT_pair_cases, RBM.sfT_pair_le}\leanok
  \uses{def:sfT}
  For $|x-\alpha|\vee|y-\alpha|\le\ell$,
  $\mathsf T_t(|x-\alpha|\wedge\ell)\mathsf T_t(|y-\alpha|\wedge\ell)\lesssim
   \mathsf T_t(|x-y|\wedge\ell)\,\Psi_t\,w_d(|x-\alpha|\wedge|y-\alpha|)$, and when
  $|y-\alpha|>\ell$, $\mathsf T_t(|x-\alpha|\wedge\ell)\mathsf T_t(\ell)\le
   \mathsf T_t(\ell)\Psi_t\,w_d(|x-\alpha|\wedge\ell)$ with constant $1$.
  Every pair $(x,y)$ and every $\alpha$ falls under one of the two (\texttt{sfT\_pair\_cases}),
  which is what makes the index ranges of the third proofreading round -- $\mathtt{(eq{:}KtKt)}$
  for $2\le i\le k$ in case 2 and for $1\le i\le k$ in case 3 -- exactly right.

  \textbf{Sharper than the paper:} the three cases merge into the single bound
  \texttt{sfT\_pair\_le}, whose polynomial factor is indexed by
  $|x-\alpha|\wedge|y-\alpha|\wedge\ell$.  So the $2^{2k}$ regions $\mathbb D_{\le\ell,\bm\sigma}$
  and the three-case split are not needed in Lean; the case distinction only decides
  \emph{how many} polynomial factors are kept.
\end{lemma}

\begin{lemma}[The lattice sum of Appendix A.4]\label{lem:ballsum}
  \lean{RBM.sum_shift, RBM.ballC, RBM.sum_ball_min_pow_le}\leanok
  \uses{def:lattice}
  For $d=k+2$, $\ell\ge1$ and any finite $D$ inside the ball of radius $\ell$ around $a$,
  $\sum_{\alpha\in D}(|x-\alpha|\wedge\ell+1)^{-(d-2)}\le C_d\,\ell^2$, uniformly in $x$ and $L$.

  \textbf{Sharper than the paper:} no volume count.  The paper bounds the region
  $|x-\alpha|>\ell$, where the summand is the constant $(\ell+1)^{-(d-2)}$, by
  $|\mathbb D_{\le\ell}|\lesssim\ell^d$ times that constant.  Here $\alpha\in D$ already gives
  $|a-\alpha|\le\ell$, so the K2 summand \emph{centred at $a$} is itself
  $\ge e^{-1}(\ell+1)^{-(d-2)}$: both regions are paid for by \texttt{sum\_radial\_exp\_le},
  and the $\ell^2$ is the same $\ell^2$.
\end{lemma}

\begin{lemma}[$\mathtt{claim{:}TTk}$, $\mathtt{(eq{:}key\_T\_reudce)}$ summed]\label{lem:TTk-sum}
  \lean{RBM.prod_sfT_pair_le, RBM.prod_wfac_le_two, RBM.keyC, RBM.key_T_reduce}\leanok
  \uses{lem:TTk-pointwise, lem:ballsum}
  For $k\ge2$, $\ell\ge1$ and $D$ inside the ball of radius $\ell$ around $a$,
  \[\sum_{\alpha\in D}\prod_{i=1}^k\mathsf T_t(|x_i-\alpha|\wedge\ell)\mathsf T_t(|y_i-\alpha|\wedge\ell)
    \;\le\;C(d,k)\,\bigl(\Psi_t^2\ell^2\bigr)\,\Psi_t^{k-2}\prod_{i=1}^k\mathsf T_t(|x_i-y_i|\wedge\ell).\]
  The hypothesis $k\ge2$ is essential -- the paper notes that the estimate fails for $k=1$ --
  and is what lets \texttt{prod\_wfac\_le\_two} keep two polynomial factors.

  \textbf{Not formalized yet (Q29):} the passage from $\Psi_t^2\ell^2$ to $(W^d\eta_t)^{-1}$,
  which uses $\ell\le(\log W)^{10}\ell_t$ and $\ell_t^2B_{t,0}\lesssim|1-t|^{-1}$ under
  $1-t\ge\hat{\gcoup}^2/L^2$.  That is where the logarithms enter and why the paper states
  the result with $\prec$ rather than $\lesssim$.
\end{lemma}

% ---------------------------------------------------------------------------
% lem:sum_decay
% ---------------------------------------------------------------------------

\begin{lemma}[$\mathtt{(eq{:}latticesum\_d3)}$, the estimate the third round added]\label{lem:latticesum}
  \lean{RBM.sum_inv_Icc_le, RBM.sum_radial_tail_le, RBM.latC, RBM.latticesum_d3}\leanok
  \uses{def:lattice, def:params}
  $\sum_{|x|>\ell}(|x|+1)^{-(d-1)}e^{-\sqrt{|x|/\ell}}\lesssim \ell\,\log\ell$ and the
  lattice sum $\mathtt{(eq{:}latticesum\_d3)}$ itself.  Proved for every $d\ge3$ at once,
  without splitting $d=3$ from $d\ge4$; it uses no interface hypothesis.
\end{lemma}

\begin{lemma}[$\mathtt{(eq{:}decomp\_U2)}$ and $\mathtt{(eq{:}decayXi)}$]\label{lem:sumdecay-parts}
  \lean{RBM.XiKer, RBM.uKer_eq_one_add_XiKer, RBM.SB_apply_eq_zero_of_one_lt, RBM.UN_apply_eq_sum_powerset, RBM.norm_XiKer_apply_le}\leanok
  \uses{lem:sum-Ndecay, bor:ThfadC}
  $\Ucal^{(n)}$ expands over subsets ($\prod(1+\Xi^{(i)})$, via \texttt{Finset.prod\_add}), and
  the entries of $\Xi$ decay: the zero-mode term is absorbed by
  \texttt{zeroMode\_le\_of\_ge\_mul}, the nearest-neighbour support of $\SB$ moves the profile
  from $|c-b|$ to $|a-b|$, and the row sum $1$ eats the sum over $c$.

  \textbf{Not formalized yet (Q28):} the three conclusions of $\mathtt{lem{:}sum\_decay}$
  themselves ($\mathtt{sum\_res\_1}$, the non-alternating case, the sum-zero property).
  The trap recorded there: going through $\|\Xi\|_{\infty\to\infty}$ gives a shape weaker
  than the claim; one has to use the ball truncation of $\mathtt{(deccA0)}$.
\end{lemma}

% ---------------------------------------------------------------------------
% Appendix A.5: K-loops
% ---------------------------------------------------------------------------

\begin{definition}[Canonical partitions of a polygon, $\mathtt{TSP}(\mathcal P_a)$]\label{def:tsp}
  \lean{RBM.Loop.diagonals, RBM.Loop.CrossingFree, RBM.Loop.TSP, RBM.Loop.thetaEdge, RBM.Loop.polyVal, RBM.Loop.treeVal, RBM.Loop.treeSum, RBM.Loop.GammaSum}\leanok
  \uses{def:Theta}
  Crossing-free sets of diagonals of an oriented polygon, the value $\Gamma^{(n)}$ of a
  canonical partition built from $\Th$-edges by the recursion \texttt{polyVal}, and the
  sums $\sum_{\Gamma\in\mathrm{TSP}}$.  The representation is the same as in the sister
  project \texttt{RBM1D}: it is dimension-independent.
\end{definition}

\begin{theorem}[$\mathtt{ML{:}Kbound}$, and the $d\ge3$ step --- \emph{statement assumed, the modification proved}]\label{hyp:Kbound}
  \lean{RBM.Loop.KLoopBound, RBM.Loop.inv_pow_pair_le, RBM.Loop.not_inv_pow_pair_le_single}
  \uses{def:params, def:kloop, bor:ThfadC, bor:ThfadCshort, lem:pureloop}
  $\max_{\bm\sigma}\max_{\mathbf a}|\Kcal^{(n)}_{t,\bm\sigma,\mathbf a}|\prec(W^{-d}B_{t,0})^{n-1}$
  for every $n$ and $t\in[0,1)$.  Appendix A.5 says the proof is analogous to
  \texttt{[YY\_25]} Lemma 3.11 \emph{but requires additional modifications to handle
  $d\ge3$}, without saying which.

  \textbf{It is this step}, in case (ii).4 of $\mathtt{(eq{:}ind\text{-}step\text{-}bound)}$:
  \[(x^{d-1}+1)^{-1}(y^{d-1}+1)^{-1}\le 2\bigl[(x^{d}+1)^{-1}(y^{d-2}+1)^{-1}+(x^{d-2}+1)^{-1}(y^{d}+1)^{-1}\bigr],\]
  proved as \texttt{inv\_pow\_pair\_le}.  The asymmetry is forced: bounding one
  $(d-1)$-factor by $1$ leaves a divergent lattice sum, and even
  $\sum_b(|a-b|^d+1)^{-1}$ is of order $\log L$; what makes the sum converge is the
  leftover exponent $d-2$, which is useful exactly when $d\ge3$.  In $d=1$,
  $|\cdot|^{d-1}=1$ carries no decay at all, which is why the argument of \texttt{[YY\_25]}
  has to be rewritten.  \texttt{not\_inv\_pow\_pair\_le\_single} checks that one term of
  the split does not suffice.

  The statement itself is carried as a hypothesis: its proof needs the molecule layer
  $\Sigma^{(\pi)}$ and the sum-zero property, which are not formalized.  It is
  \emph{owed}, not borrowed: Appendix A.5 does give the proof (``the proof is analogous to
  that of Lemma 3.11 in \texttt{[YY\_25]} \dots for the reader's convenience, we provide
  the proof below''), so this is a debt of the Lean development.  The remaining
  half of the $d\ge3$ step, the lattice sum
  $\sum_b(|a-b|^d+1)^{-1}(|c-b|^{d-2}+1)^{-1}\lesssim1$, is work order Q32.
\end{theorem}

\begin{lemma}[$\mathtt{(eq\_Ktree)}$ at $n=3$]\label{lem:treeThree}
  \lean{RBM.Loop.treeEqRhs_three, RBM.Loop.treeSum_three, RBM.Loop.sum_SB_starLeft, RBM.Loop.sum_SB_starRight, RBM.Loop.hasDerivAt_starThree, RBM.Loop.kThree, RBM.Loop.hasDerivAt_kThree, RBM.Loop.kThree_zero, RBM.Loop.kThree_eq_of_isKLoop}\leanok
  \uses{def:tsp, def:kloop, hyp:twoloop, lem:unique, lem:KTwoFormula, lem:Theta-deriv}
  A triangle has no diagonals, so the tree sum is the single star
  $\sum_b\Th^{(\sigma_0\sigma_1)}(a_0,b)\Th^{(\sigma_1\sigma_2)}(a_1,b)\Th^{(\sigma_2\sigma_0)}(a_2,b)$,
  and it is the $3$-loop of \emph{every} family of $\Kcal$-loops whose $2$-loops are
  bounded on each $[0,T_0]$, $T_0<1$.

  The content is the match between the three boundary edges and the three $(k,l)$ terms:
  $(1,2)$ cuts off the $2$-chain $(\sigma_0,\sigma_1)$ and replaces $a_0$, $(2,3)$ cuts off
  $(\sigma_1,\sigma_2)$ and replaces $a_1$, and $(1,3)$ cuts off $(\sigma_0,\sigma_2)$ and
  replaces $a_2$ --- that last one with its $2$-chain on the \emph{left} of $\SB$, which is
  why there are two summation lemmas and not one.  Differentiating the edge at $a_i$
  inserts $\SB$ there, and the matching $(k,l)$ term rebuilds exactly that sandwich.

  Only the $2$-loop bound is needed, never a bound on $3$-loops: for $n\ge3$ the equation
  is linear in the $n$-loops with the $2$-loops as coefficients.
  Work orders: $n=4$, where internal edges first appear, is Q30; general $n$ is Q31.
\end{lemma}

\begin{definition}[$\Kcal$-loops: the index, the cut-and-glue, $\mathtt{(pro\_dyncalK)}$]\label{def:kloop}
  \lean{RBM.Loop.LoopIdx, RBM.Loop.LoopIdx.cutGlueL, RBM.Loop.LoopIdx.cutGlueR, RBM.Loop.treeEqRhs, RBM.Loop.treeEqRhs_two, RBM.Loop.MLoop, RBM.Loop.IsKLoop}\leanok
  \uses{def:Theta, def:tsp}
  A loop index is a pair of lists $(\bm\sigma,\mathbf a)$ of equal length; cutting the
  polygon at $(k,l)$ and gluing in an internal vertex $b$ gives the two shorter indices
  \texttt{cutGlueL}/\texttt{cutGlueR}, whose lengths add up to $n+2$
  (\texttt{length\_cutGlueL\_add\_length\_cutGlueR}).  \texttt{treeEqRhs} is the right-hand
  side of $\mathtt{(pro\_dyncalK)}$ assembled from them, \texttt{MLoop} is the value at
  $t=0$, and \texttt{IsKLoop} says that a family $K$ satisfies the equation with that
  initial condition.

  This node is \emph{definitions}, not a result: it says what the objects are.  The audit
  keeps \texttt{IsKLoop} in its structural list for exactly that reason --- assuming it is
  saying what the data is, not borrowing a theorem.  \texttt{treeEqRhs\_two} is the
  evaluation at $n=2$ that the $2$-loop work runs on.
\end{definition}

\begin{theorem}[Tree representation of $\Kcal$-loops --- \emph{borrowed}]\label{bor:KTreeRep}
  \lean{RBM.Loop.KTreeRep}
  \uses{def:kloop, def:tsp}
  $\Kcal^{(n)}_{t,\bm\sigma,\mathbf a}=W^{-d(n-1)}\sum_{\Gamma\in\mathrm{TSP}(\mathcal P_{\mathbf a})}\Gamma^{(n)}$.
  \emph{Attributed to} \texttt{[YY\_25]} Lemma 3.4 --- the paper recalls it, and does not
  prove it, which is why this node is \texttt{bor:}.  The statement layer is in Lean (Q15);
  uniqueness is no longer assumed (\ref{lem:unique}, Q22a) and the cases $n=2$
  (\ref{lem:KTwoFormula}) and $n=3$ (\ref{lem:treeThree}) are proved, so what is still
  carried as a hypothesis is existence for $n\ge4$: Q30 ($n=4$, internal edges) and
  Q31 (general $n$).  Those two work orders would leave the \texttt{bor:} ledger one
  entry shorter than the paper's own bibliography --- a case where the formalization
  proves more than the paper does.
\end{theorem}

\begin{theorem}[The a priori bound on $2$-loops --- \emph{owed}]\label{hyp:twoloop}
  \lean{RBM.Loop.TwoLoopBounded}
  \uses{def:kloop}
  On each $[0,T_0]$ with $T_0<1$ the $2$-loops of $K$ are bounded.  Every uniqueness
  statement here takes this as a hypothesis (\ref{lem:unique}, \ref{lem:treeThree}).

  It is \emph{owed}, not borrowed: for the explicit solution the bound is already proved
  --- \texttt{norm\_kTwo\_le} of \ref{lem:kTwo} gives $W^{-d}(1-t)^{-1}$ --- and the paper
  obtains it along the way rather than citing it.  What is missing is the a priori form,
  for a family of $\Kcal$-loops not yet known to be that solution.
\end{theorem}

\begin{lemma}[Uniqueness for $\mathtt{(pro\_dyncalK)}$]\label{lem:unique}
  \lean{RBM.Loop.LoopIdx.length_cutGlueR_eq_two, RBM.Loop.LoopIdx.length_cutGlueL_eq_two, RBM.Loop.eq_on_level, RBM.Loop.isKLoop_unique}\leanok
  \uses{def:kloop, hyp:twoloop}
  Two families of $\Kcal$-loops on $[0,T_0]$ with the same $W$, $g$, $m$, whose $2$-loops
  stay bounded, agree on every loop of length $\ge2$.

  One Grönwall argument covers all lengths.  The structure of $\mathtt{(pro\_dyncalK)}$ is
  what makes it work: the two chains produced by a cut have lengths adding up to $n+2$, so
  if one of them has the full length $n$ the other is a $2$-loop.  Hence for $n\ge3$ the
  equation at length $n$ is \emph{linear} in the length-$n$ unknowns once the shorter
  lengths are fixed, with the $2$-loops as coefficients, and for $n=2$ it is a Riccati
  equation; the same estimate covers both, and strong induction on $n$ finishes.
\end{lemma}

\begin{theorem}[$\mathtt{(Kn2sol)}$, the $n=2$ loop --- \emph{now a theorem}]\label{lem:KTwoFormula}
  \lean{RBM.Loop.KTwoFormula, RBM.Loop.kTwoFormula_of_isKLoop, RBM.Loop.pureLoop_two_of_isKLoop}\leanok
  \uses{lem:kTwo, lem:unique}
  The closed formula for $\Kcal^{(2)}$.  The paper attributes it to \texttt{[YY\_25]},
  \texttt{[RBSO1D]}; it is not an input.  \texttt{kTwoFormula\_of\_isKLoop} proves it of
  \emph{every} family of $\Kcal$-loops on $[0,1)$ whose $2$-loops are bounded on each
  $[0,T_0]$, $T_0<1$: existence (\ref{lem:kTwo}) plus uniqueness at length $2$
  (\ref{lem:unique}, where the hypothesis on shorter loops is vacuous).

  What is left is the a priori bound on the $2$-loops, which the paper proves along the
  way rather than borrows.  \texttt{pureLoop\_two\_of\_isKLoop} is \ref{lem:pureloop} with
  this hypothesis gone, resting only on \ref{bor:ThfadCshort}.
\end{theorem}

\begin{lemma}[$\mathtt{(Kn2sol)}$, the two-loop solution --- \emph{proved, not borrowed}]\label{lem:kTwo}
  \lean{RBM.Loop.kTwo, RBM.Loop.hasDerivAt_kTwo, RBM.Loop.kTwo_zero, RBM.Loop.kTwoLoop, RBM.Loop.hasDerivAt_kTwoLoop, RBM.Loop.kTwoFormula_kTwoLoop, RBM.Loop.norm_kTwo_le, RBM.Loop.pureLoop_two_kTwoLoop}\leanok
  \uses{lem:Theta-deriv, def:tsp, def:kloop}
  $K^{(2)}_{t,\bm\sigma,\mathbf a}=W^{-d}m(\sigma_1)m(\sigma_2)\,(\Th_{t\,m(\sigma_1)m(\sigma_2)})_{a_1a_2}$
  solves the convolution tree equation at $n=2$ and takes the $M$-loop value at $t=0$.
  The paper records this in an example and attributes it to \texttt{[YY\_25]},
  \texttt{[RBSO1D]}; it does not have to be borrowed, and here it is not.

  This is the \emph{existence} half; together with uniqueness (\ref{lem:unique}) it gives
  \ref{lem:KTwoFormula} for every family of $\Kcal$-loops.  \texttt{norm\_kTwo\_le} is the
  entry bound $W^{-d}(1-t)^{-1}$ for the explicit solution, which is what supplies the
  $2$-loop bound the uniqueness argument asks for.
\end{lemma}

\begin{lemma}[$\mathtt{lem\_pureloop}$, $n=2$]\label{lem:pureloop}
  \lean{RBM.Loop.norm_Theta_same_le_exp, RBM.Loop.sum_exp_decay_conv, RBM.Loop.pureLoop_two}\leanok
  \uses{lem:KTwoFormula, bor:ThfadCshort}
  For a pure loop ($\sigma_1=\sigma_2$), $|\Kcal^{(2)}_{t,\bm\sigma,\mathbf a}|\le
  C W^{-d}e^{-c|a_1-a_2|}$.  The two tools are the exponential decay of $\Th^{(\sigma,\sigma)}$
  (from $\mathtt{(prop{:}ThfadC\_short)}$) and a convolution of exponentials, in which the
  decay constant halves at each step.

  \textbf{Not formalized yet (Q25):} general $n$, by induction on the \texttt{polyVal}
  recursion, with the induction hypothesis stated for the \emph{maximal pairwise} distance
  of the label set.
\end{lemma}
